% Preámbulo común para las soluciones de los prácticos de Programación 3
\documentclass[11pt,a4paper]{article}

\usepackage[utf8]{inputenc}
\usepackage[T1]{fontenc}
\usepackage{lmodern}
\usepackage[spanish,es-noshorthands]{babel}
\usepackage[margin=2.3cm]{geometry}
\usepackage{amsmath,amssymb,amsthm}
\usepackage{enumitem}
\usepackage{xcolor}
\usepackage{tikz}
\usetikzlibrary{arrows.meta,positioning,shapes,calc}
\usepackage[noend]{algpseudocode}
\usepackage{algorithm}
\usepackage{booktabs}
\usepackage{tcolorbox}
\usepackage[hidelinks]{hyperref}

\floatname{algorithm}{Algoritmo}
\renewcommand{\algorithmicrequire}{\textbf{Entrada:}}
\renewcommand{\algorithmicensure}{\textbf{Salida:}}
\algrenewcommand\algorithmicwhile{\textbf{mientras}}
\algrenewcommand\algorithmicdo{\textbf{hacer}}
\algrenewcommand\algorithmicfor{\textbf{para}}
\algrenewcommand\algorithmicforall{\textbf{para todo}}
\algrenewcommand\algorithmicif{\textbf{si}}
\algrenewcommand\algorithmicthen{\textbf{entonces}}
\algrenewcommand\algorithmicelse{\textbf{sino}}
\algrenewcommand\algorithmicreturn{\textbf{retornar}}
\algrenewcommand\algorithmicfunction{\textbf{función}}
\algrenewcommand\algorithmicprocedure{\textbf{procedimiento}}

\theoremstyle{plain}
\newtheorem{lema}{Lema}[section]
\newtheorem{teorema}[lema]{Teorema}
\newtheorem{proposicion}[lema]{Proposición}
\theoremstyle{definition}
\newtheorem{definicion}[lema]{Definición}
\theoremstyle{remark}
\newtheorem*{observacion}{Observación}

% Caja para el enunciado de cada ejercicio
\newtcolorbox{enunciado}{colback=gray!6,colframe=gray!50,boxrule=0.4pt,arc=2pt,left=6pt,right=6pt,top=4pt,bottom=4pt,fontupper=\small}
% Nota que vincula el ejercicio con parciales/exámenes
\newcommand{\examref}[1]{\par\noindent{\small\color{blue!50!black}\textbf{Relación con parciales/exámenes:} #1}\par\medskip}

\newcommand{\NP}{\mathsf{NP}}
\newcommand{\Pclass}{\mathsf{P}}
\newcommand{\Oh}{\mathcal{O}}

\setlength{\parskip}{0.4em}
\setlength{\parindent}{0pt}


\title{Práctico 2 --- Análisis de algoritmos}
\author{Programación 3 -- FING, UdelaR}
\date{2026}

\begin{document}
\maketitle

\paragraph{Definiciones.} Para funciones $T,f:\mathbb{N}\to\mathbb{R}_{\ge 0}$:
\begin{itemize}
  \item $T(n)$ es $\Oh(f(n))$ si existen $c>0$ y $n_0\ge 0$ tales que $T(n)\le c\,f(n)$ para
        todo $n\ge n_0$.
  \item $T(n)$ es $\Omega(f(n))$ si existen $\varepsilon>0$ y $n_0\ge 0$ tales que
        $T(n)\ge \varepsilon\, f(n)$ para todo $n\ge n_0$.
  \item $T(n)$ es $\Theta(f(n))$ si es $\Oh(f(n))$ y $\Omega(f(n))$.
\end{itemize}

%=====================================================================
\section{Ejercicio 1}

\begin{enunciado}
(K\&T, Ex.~2.5) Sean $f$ y $g$ funciones positivas tales que $f(n)$ es $\Oh(g(n))$. Para cada
una de las siguientes afirmaciones, indique si es verdadera, demostrándolo, o falsa, dando un
contraejemplo.
\textbf{(a)} $f(n)^2$ es $\Oh(g(n)^2)$.
\textbf{(b)} $2^{f(n)}$ es $\Oh(2^{g(n)})$.
\end{enunciado}

\examref{Examen diciembre 2017, Pregunta~2 (a) y (b) son exactamente estas dos afirmaciones;
la parte (b) también aparece en el Examen julio 2017, Pregunta~1(b), y en el Primer parcial 2016,
Ej.~1(a)\,I (con $g(n)=n$). Las soluciones oficiales dan la misma demostración para (a) y
contraejemplos del mismo tipo para (b) ($f(n)=2n$, $g(n)=n$).}

\subsection*{Parte (a): verdadera}
Como $f(n)$ es $\Oh(g(n))$, existen $c>0$ y $n_0$ tales que $f(n)\le c\,g(n)$ para todo
$n\ge n_0$. Como $f$ y $g$ son positivas, ambos lados de la desigualdad son positivos y
elevar al cuadrado (función creciente en $\mathbb{R}_{\ge0}$) la preserva:
\[
f(n)^2\le c^2\,g(n)^2\qquad\text{para todo } n\ge n_0 .
\]
Tomando la constante $c'=c^2>0$ y el mismo $n_0$, se cumple la definición: $f(n)^2$ es
$\Oh(g(n)^2)$.

\subsection*{Parte (b): falsa}
Contraejemplo: $f(n)=2n$, $g(n)=n$. Claramente $f(n)$ es $\Oh(g(n))$ (con $c=2$, $n_0=0$).
Sin embargo $2^{f(n)}=4^n$ y $2^{g(n)}=2^n$. Supongamos por absurdo que existen $c>0$ y $n_0$
tales que $4^n\le c\,2^n$ para todo $n\ge n_0$. Dividiendo entre $2^n>0$ queda $2^n\le c$
para todo $n\ge n_0$, lo cual es falso para cualquier $n>\max\{n_0,\log_2 c\}$. Por lo tanto
$2^{f(n)}$ no es $\Oh(2^{g(n)})$.

(Equivalentemente, por la regla del límite: $\lim_{n\to\infty} 4^n/2^n=\lim_{n\to\infty}2^n=\infty$.)
El problema es que una constante multiplicativa en el exponente se convierte en un cambio de
base: $2^{c\,g(n)}=(2^c)^{g(n)}$.

%=====================================================================
\section{Ejercicio 2}

\begin{enunciado}
En el libro (K\&T, Cap.~2) se describe una implementación del algoritmo de Gale--Shapley cuyo
tiempo de ejecución es $\Oh(n^2)$.
\textbf{(a)} Muestre, a través de un ejemplo, que este tiempo de ejecución es también
$\Omega(n^2)$ en el peor caso, incluso si no se tiene en cuenta el tiempo de preprocesamiento en
el que se construye la matriz \texttt{Ranking}.
\textbf{(b)} El hecho que acaba de demostrar en la parte (a), ¿depende de la implementación
específica del algoritmo de Gale--Shapley que consideremos? Justifique.
\end{enunciado}

\examref{Relacionados: Examen diciembre 2021, Ej.~1 y Examen febrero 2023, Ej.~2 (la cantidad
de iteraciones de Gale--Shapley es a lo sumo $n^2$, porque cada iteración es una propuesta
nueva); Examen diciembre 2025 y Examen julio 2026, Ej.~1(a) (si algún $m$ debe proponerse a
todas las $w$).}

\paragraph{Recordatorio de la implementación del libro.} Los hombres libres se mantienen en una
lista; \texttt{Next}$[m]$ indica la próxima mujer de la lista de $m$; \texttt{Current}$[w]$ la
pareja actual de $w$; \texttt{Ranking}$[w][m]$ la posición de $m$ en la lista de $w$. Con
esto cada iteración del ciclo (una propuesta) se ejecuta en tiempo $O(1)$; como hay a lo sumo
$n^2$ propuestas, el tiempo total es $\Oh(n^2)$.

\subsection*{Parte (a)}
Para cada $n\ge1$ consideramos la instancia $I_n$ en la que \textbf{todos} los hombres tienen
la misma lista y \textbf{todas} las mujeres tienen la misma lista:
\[
m_i:\ w_1\succ w_2\succ\cdots\succ w_n,\qquad
w_j:\ m_1\succ m_2\succ\cdots\succ m_n,\qquad 1\le i,j\le n .
\]

\begin{lema}
En $I_n$ el único emparejamiento estable es $S=\{(m_i,w_i):1\le i\le n\}$.
\end{lema}
\begin{proof}
Mostramos por inducción en $k$ que todo emparejamiento estable $S'$ contiene a $(m_k,w_k)$.
Supongamos que contiene $(m_i,w_i)$ para todo $i<k$. Si $(m_k,w_k)\notin S'$, entonces $m_k$
está con alguna $w_j$ con $j>k$ (las $w_i$, $i<k$, están ocupadas) y $w_k$ está con algún
$m_l$ con $l>k$. Pero $m_k$ prefiere $w_k$ a $w_j$ (porque $k<j$) y $w_k$ prefiere $m_k$ a
$m_l$ (porque $k<l$): $(m_k,w_k)$ es una inestabilidad, absurdo. (El caso $k=1$ es
el Ejercicio~1(a) del Práctico~1.)
\end{proof}

Como Gale--Shapley devuelve un emparejamiento estable, en \emph{cualquier} ejecución sobre
$I_n$ (sea cual sea el orden en que se eligen los hombres libres) cada $m_i$ termina con $w_i$.
Ahora bien, cada hombre se propone en orden de su lista y su pareja final es la última mujer a la
que se propuso (una vez emparejado sólo vuelve a proponerse si lo abandonan). Como $w_i$ ocupa
la posición $i$ en la lista de $m_i$, el hombre $m_i$ hace exactamente $i$ propuestas:
a $w_1,\dots,w_i$. Por lo tanto la cantidad total de propuestas, es decir, de iteraciones del
ciclo principal, es
\[
\sum_{i=1}^{n} i=\frac{n(n+1)}{2}\ \ge\ \frac{n^2}{2}.
\]
Cada iteración ejecuta al menos una operación elemental (tomar un hombre libre, leer
\texttt{Next}, etc.), es decir, tiene costo al menos una constante $c>0$. Así, el tiempo de
ejecución en $I_n$, aun sin contar la construcción de \texttt{Ranking}, es al menos
$c\,n^2/2$. Concluimos que el tiempo de ejecución en el peor caso es $\Omega(n^2)$ y, junto con
la cota superior, $\Theta(n^2)$.

\paragraph{Ejemplo de ejecución, $n=3$.} Eligiendo siempre el primer hombre libre en orden de
índice: $m_1\to w_1$ (acepta); $m_2\to w_1$ (rechaza); $m_2\to w_2$ (acepta); $m_3\to w_1$,
$m_3\to w_2$ (rechazan); $m_3\to w_3$ (acepta). Son $1+2+3=6$ propuestas.

\subsection*{Parte (b)}
\textbf{No depende} de la implementación. El argumento de la parte (a) sólo usa propiedades del
\emph{algoritmo} de Gale--Shapley (cada iteración es una propuesta, cada hombre se propone en
orden de su lista y nunca dos veces a la misma mujer, el resultado es estable), y no las
estructuras de datos usadas. En particular, en $I_n$ \emph{toda} ejecución del algoritmo, con
cualquier criterio de elección del hombre libre y cualquier representación de las listas,
realiza exactamente $n(n+1)/2$ propuestas, y cada propuesta cuesta al menos tiempo constante en
cualquier implementación. Por lo tanto cualquier implementación de Gale--Shapley tiene tiempo de
peor caso $\Omega(n^2)$. Lo que sí depende de la implementación es la cota superior: una
implementación descuidada (por ejemplo, recorrer la lista de $w$ para comparar dos hombres, en
lugar de usar \texttt{Ranking}) puede costar $\Theta(n)$ por iteración y llegar a $\Theta(n^3)$.

Notar además que, si se cuenta la lectura de la entrada, cualquier algoritmo que deba leer las
$2n$ listas de largo $n$ ya insume $\Omega(n^2)$; la parte (a) muestra algo más fuerte: el ciclo
principal de Gale--Shapley por sí solo es $\Omega(n^2)$.

%=====================================================================
\section{Ejercicio 3}

\begin{enunciado}
$T(n)$ es $o(f(n))$ si para \textbf{cualquier} constante $\epsilon>0$ existe una constante
$n_\epsilon\ge 0$ tal que para todo $n\ge n_\epsilon$ se cumple $T(n)<\epsilon\cdot f(n)$.
\textbf{(a)} Demuestre que si $T(n)$ es $o(f(n))$ entonces $T(n)$ es $\Oh(f(n))$.
\textbf{(b)} Muestre un ejemplo de funciones $f,g$ para las que se cumple $f(n)$ es
$\Oh(g(n))$ y $f(n)$ no es $o(g(n))$.
\textbf{(c)} Demuestre que si $T(n)$ es $o(f(n))$ entonces $T(n)$ no es $\Omega(f(n))$.
\end{enunciado}

\subsection*{Parte (a)}
Supongamos que $T(n)$ es $o(f(n))$. Aplicando la definición con $\epsilon=1$, existe $n_1$ tal
que $T(n)<1\cdot f(n)$ para todo $n\ge n_1$. En particular $T(n)\le c\,f(n)$ para todo
$n\ge n_0$, con $c=1$ y $n_0=n_1$. Luego $T(n)$ es $\Oh(f(n))$.

\subsection*{Parte (b)}
Sean $f(n)=g(n)=n$. Se cumple $f(n)$ es $\Oh(g(n))$ con $c=1$, $n_0=0$, pues $n\le 1\cdot n$.
Pero $f(n)$ no es $o(g(n))$: para $\epsilon=1$ la definición pediría un $n_1$ tal que
$n<1\cdot n$ para todo $n\ge n_1$, y la desigualdad $n<n$ es falsa para todo $n$. (Más en
general, ninguna función positiva es $o$ de sí misma.)

\subsection*{Parte (c)}
Supongamos, por absurdo, que $T(n)$ es $o(f(n))$ y también $\Omega(f(n))$. Por ser
$\Omega(f(n))$ existen $\varepsilon>0$ y $n_0$ tales que
\[
T(n)\ge \varepsilon\, f(n)\qquad \text{para todo } n\ge n_0 .
\]
Por ser $o(f(n))$, aplicando la definición con $\epsilon=\varepsilon$ (que es una constante
positiva), existe $n_\varepsilon$ tal que
\[
T(n)< \varepsilon\, f(n)\qquad \text{para todo } n\ge n_\varepsilon .
\]
Tomando $n=\max\{n_0,n_\varepsilon\}$ se obtiene $\varepsilon f(n)\le T(n)<\varepsilon f(n)$,
una contradicción. Por lo tanto $T(n)$ no es $\Omega(f(n))$.

En particular, las partes (a) y (c) dicen que $o(f)$ consiste en cotas superiores \emph{no
ajustadas}: $T\in o(f)$ implica $T\in \Oh(f)\setminus \Theta(f)$.

%=====================================================================
\section{Ejercicio 4}

\begin{enunciado}
Dada una secuencia $S=s_1\dots s_n$ de enteros con valores en el rango $1\dots m$, donde $m$
está determinado por una función conocida de $n$, el siguiente algoritmo devuelve los elementos
de $S$ en orden creciente. Analice su tiempo de ejecución.
\begin{algorithmic}
\Statex \textbf{Algoritmo} \textsc{Ordenar}$_{1\_m}$ \quad (Entrada: $S=s_1\dots s_n$; Resultado: $R$)
\State Crear un arreglo $C$ de contadores con índices de $1$ a $m$, inicializados en $0$
\For{$i$ desde $1$ hasta $n$}
  \State Sumar 1 a $C[s_i]$
\EndFor
\For{$i$ desde $1$ hasta $m$}
  \For{$j$ desde $1$ hasta $C[i]$}
    \State Agregar $i$ al final de $R$
  \EndFor
\EndFor
\end{algorithmic}
\end{enunciado}

(Siguiendo la numeración de líneas del enunciado original: creación de $C$ en la línea 2, primer ciclo
en las líneas 3--4 y ciclos anidados en las líneas 5--7.)

\paragraph{Hipótesis de costo.} Acceder o modificar una celda de un arreglo es $O(1)$, y
agregar un elemento al final de $R$ es $O(1)$ (por ejemplo, $R$ es una lista con puntero al
último elemento, o un arreglo de tamaño $n$ con un índice de la próxima posición libre).

\paragraph{Línea 2.} Crear e inicializar $m$ contadores requiere tiempo $\Theta(m)$.

\paragraph{Líneas 3--4.} El ciclo itera exactamente $n$ veces y cada iteración hace un acceso a
$s_i$ y un incremento de $C[s_i]$, en tiempo $O(1)$ (el índice $s_i$ está en $1\dots m$, así
que el acceso es válido). Total: $\Theta(n)$.

Al terminar este ciclo, $C[i]$ es la cantidad de apariciones de $i$ en $S$, y por lo tanto
\[
\sum_{i=1}^{m} C[i]=n .
\]

\paragraph{Líneas 5--7.} El ciclo externo itera $m$ veces. En la iteración $i$, el control del
ciclo externo y la inicialización del interno cuestan $O(1)$, y el ciclo interno itera $C[i]$
veces con cuerpo $O(1)$. Entonces el costo es, para ciertas constantes $a,b>0$,
\[
\sum_{i=1}^{m}\bigl(a+b\,C[i]\bigr)=a\,m+b\sum_{i=1}^{m}C[i]=a\,m+b\,n,
\]
es decir $\Theta(m+n)$. Notar que el ciclo interno \emph{no} cuesta $\Oh(n)$ por cada $i$
(lo que daría la cota pesimista $\Oh(mn)$): el total de iteraciones internas, sumando sobre todos
los $i$, es exactamente $n$.

\paragraph{Total.} Sumando, el tiempo de ejecución es
\[
T(n)=\Theta(m)+\Theta(n)+\Theta(m+n)=\Theta(n+m(n)).
\]
Como cada línea se ejecuta la misma cantidad de veces para cualquier entrada de tamaño $n$
(el costo no depende de los valores de $S$), este es a la vez el tiempo de mejor caso, de peor
caso y de caso promedio. Según cómo crece $m$ en función de $n$:
\begin{itemize}
  \item si $m(n)$ es $\Oh(n)$ (por ejemplo $m=n$ o $m$ constante), el algoritmo es
        $\Theta(n)$: \textbf{lineal}, mejor que la cota $\Omega(n\log n)$ de los ordenamientos
        por comparaciones, cota que no aplica porque este algoritmo no compara elementos sino que
        los usa como índices;
  \item si $m(n)$ crece más rápido que $n$ (por ejemplo $m=n^2$ o $m=2^n$), domina el término
        $m$ y el algoritmo es $\Theta(m(n))$, que puede ser peor que $\Theta(n\log n)$ (ya con
        $m=n^2$) e incluso exponencial.
\end{itemize}
El espacio adicional usado es $\Theta(m)$ para $C$ más $\Theta(n)$ para $R$.

\end{document}
